\documentclass[10pt, letterpaper]{article}

% Packages:
\usepackage[
    ignoreheadfoot, % set margins without considering header and footer
    top=2 cm, % seperation between body and page edge from the top
    bottom=2 cm, % seperation between body and page edge from the bottom
    left=2 cm, % seperation between body and page edge from the left
    right=2 cm, % seperation between body and page edge from the right
    footskip=1.0 cm, % seperation between body and footer
    % showframe % for debugging 
]{geometry} % for adjusting page geometry
\usepackage[explicit]{titlesec} % for customizing section titles
\usepackage{tabularx} % for making tables with fixed width columns
\usepackage{array} % tabularx requires this
\usepackage[dvipsnames]{xcolor} % for coloring text
\definecolor{primaryColor}{RGB}{0, 79, 144} % define primary color
\usepackage{enumitem} % for customizing lists
\usepackage{fontawesome5} % for using icons
\usepackage{amsmath} % for math
\usepackage[
    pdftitle={Curriculum Fernando Gonzalez},
    pdfauthor={Fernando Gonzalez Gonzalez},
    pdfcreator={LaTeX with RenderCV},
    colorlinks=true,
    urlcolor=primaryColor
]{hyperref} % for links, metadata and bookmarks
\usepackage[pscoord]{eso-pic} % for floating text on the page
\usepackage{calc} % for calculating lengths
\usepackage{bookmark} % for bookmarks
\usepackage{lastpage} % for getting the total number of pages
\usepackage{changepage} % for one column entries (adjustwidth environment)
\usepackage{paracol} % for two and three column entries
\usepackage{ifthen} % for conditional statements
\usepackage{needspace} % for avoiding page brake right after the section title
\usepackage{iftex} % check if engine is pdflatex, xetex or luatex

% Ensure that generate pdf is machine readable/ATS parsable:
\ifPDFTeX
    \input{glyphtounicode}
    \pdfgentounicode=1
    \usepackage[T1]{fontenc}
    \usepackage[utf8]{inputenc}
    \usepackage{lmodern}
\fi

\usepackage[default, type1]{sourcesanspro} 

% Some settings:
\AtBeginEnvironment{adjustwidth}{\partopsep0pt} % remove space before adjustwidth environment
\pagestyle{empty} % no header or footer
\setcounter{secnumdepth}{0} % no section numbering
\setlength{\parindent}{0pt} % no indentation
\setlength{\topskip}{0pt} % no top skip
\setlength{\columnsep}{0.15cm} % set column seperation
\makeatletter
\let\ps@customFooterStyle\ps@plain % Copy the plain style to customFooterStyle
\patchcmd{\ps@customFooterStyle}{\thepage}{
    \color{gray}\textit{\small Fernando González González}
}{}{} % replace number by desired string
\makeatother
\pagestyle{customFooterStyle}

\titleformat{\section}{
    % avoid page braking right after the section title
    \needspace{4\baselineskip}
    % make the font size of the section title large and color it with the primary color
    \Large\color{primaryColor}
}{
}{
}{
    % print bold title, give 0.15 cm space and draw a line of 0.8 pt thickness
    % from the end of the title to the end of the body
    \textbf{#1}\hspace{0.15cm}\titlerule[0.8pt]\hspace{-0.1cm}
}[] % section title formatting

\titlespacing{\section}{
    % left space:
    -1pt
}{
    % top space:
    0.3 cm
}{
    % bottom space:
    0.2 cm
} % section title spacing

% \renewcommand\labelitemi{$\vcenter{\hbox{\small$\bullet$}}$} % custom bullet points
\newenvironment{highlights}{
    \begin{itemize}[
        topsep=0.10 cm,
        parsep=0.10 cm,
        partopsep=0pt,
        itemsep=0pt,
        leftmargin=0.4 cm + 10pt
    ]
}{
    \end{itemize}
} % new environment for highlights

\newenvironment{highlightsforbulletentries}{
    \begin{itemize}[
        topsep=0.10 cm,
        parsep=0.10 cm,
        partopsep=0pt,
        itemsep=0pt,
        leftmargin=10pt
    ]
}{
    \end{itemize}
} % new environment for highlights for bullet entries


\newenvironment{onecolentry}{
    \begin{adjustwidth}{
        0.2 cm + 0.00001 cm
    }{
        0.2 cm + 0.00001 cm
    }
}{
    \end{adjustwidth}
} % new environment for one column entries

\newenvironment{twocolentry}[2][]{
    \onecolentry
    \def\secondColumn{#2}
    \setcolumnwidth{\fill, 4.5 cm}
    \begin{paracol}{2}
}{
    \switchcolumn \raggedleft \secondColumn
    \end{paracol}
    \endonecolentry
} % new environment for two column entries

\newenvironment{threecolentry}[3][]{
    \onecolentry
    \def\thirdColumn{#3}
    \setcolumnwidth{0 cm, \fill, 4.5 cm}
    \begin{paracol}{3}
    \noindent{\raggedright #2}\switchcolumn
}{
    \switchcolumn \raggedleft \thirdColumn
    \end{paracol}
    \endonecolentry
} % new environment for three column entries

\newenvironment{header}{
    \setlength{\topsep}{0pt}\par\kern\topsep\centering\color{primaryColor}\linespread{1.5}
}{
    \par\kern\topsep
} % new environment for the header

% save the original href command in a new command:
\let\hrefWithoutArrow\href

% new command for external links:
\renewcommand{\href}[2]{\hrefWithoutArrow{#1}{\ifthenelse{\equal{#2}{}}{ }{#2 }\raisebox{.15ex}{\footnotesize \faExternalLink*}}}


\begin{document}
    \newcommand{\AND}{\unskip
        \cleaders\copy\ANDbox\hskip\wd\ANDbox
        \ignorespaces
    }
    \newsavebox\ANDbox
    \sbox\ANDbox{}

    \begin{header}
        \fontsize{30 pt}{30 pt}
        \textbf{Fernando González González}

        \vspace{0.3 cm}

        \normalsize
        \mbox{{\footnotesize\faMapMarker*}\hspace*{0.13cm}Ferrol, A Coruña}%
        \kern 0.25 cm%
        \AND%
        \kern 0.25 cm%
        \mbox{\hrefWithoutArrow{mailto:fernando.ditos@gmail.com}{{\footnotesize\faEnvelope[regular]}\hspace*{0.13cm}fernando.ditos@gmail.com}}%
        \kern 0.25 cm%
        \AND%
        \kern 0.25 cm%
        \mbox{\hrefWithoutArrow{tel:+34 671 801 801}{{\footnotesize\faPhone*}\hspace*{0.13cm}671 801 801}}%
        \kern 0.25 cm%
        \AND%
        \kern 0.25 cm%
        \mbox{\hrefWithoutArrow{https://www.linkedin.com/in/fernando-gonzalez-gonzález-7b342740b/}{{\footnotesize\faLinkedin}\hspace*{0.13cm}LinkedIn}}%
        \kern 0.25 cm%
        \AND%
        \kern 0.25 cm%
        \mbox{\hrefWithoutArrow{https://github.com/fernandogg2004}{{\footnotesize\faGithub}\hspace*{0.13cm}GitHub}}%
        \kern 0.25 cm%

    \end{header}

    \vspace{0.3 cm - 0.3 cm}


    \section{}



        
        \begin{onecolentry}
            Data Scientist · Predictive Maintenance \& Industrial IoT | Python · XGBoost · BigQuery
        \end{onecolentry}


    \section{Formación}
    
        

        \medskip
        
        \begin{threecolentry}{}{
            Sep 2022 - Jun 2026
        }
            \textbf{Grado en Matemáticas}\\
            \textit{Universidad de Santiago de Compostela} 
            \begin{highlights}
                \item \textbf{Nota media:} $8.6$.
                \item \textbf{Logros:} 8 Matrículas de Honor (incluyendo Álgebra Lineal, Geometría, Topología y Series de Fourier).
                \item \textbf{TFG:} "Modelos de regresión no lineales".
            \end{highlights}
        \end{threecolentry}

        \medskip

        \begin{threecolentry}{}{Sep 2020 - Jun 2021}
            \textbf{Academic Exchange Year}\\
            \textit{Neville High School, Luisiana, EE. UU.}
        \end{threecolentry}


    \medskip
    
    \section{Experiencia Profesional}


        \begin{threecolentry}{}{
            Jul 2025 - Ago 2025
        }
            \textbf{Data Scientist — Prácticas de Alta Intensidad} $\mid$ \textit{EcoMT}  
            \begin{highlights}
                \item \textbf{Desplegué en producción un sistema de detección temprana de fallos mecánicos con 97\% de F1-Score ponderado}, adelantando averías hasta 7 días y reduciendo potencialmente costes de mantenimiento no planificado sobre una flota de +1.000 activos industriales.
                \item  \textbf{Diseñé el pipeline completo de feature engineering sobre 1,5M de registros de series temporales en BigQuery}, incorporando rolling statistics y transformadas de Fourier (FFT) sobre señales de sensor, validado con cross-validation estratificada para evitar data leakage temporal.
                \item \textbf{Reduje un 95\% el ruido en el pipeline de datos} automatizando la detección de anomalías de sensores con Isolation Forest, eliminando entradas erróneas y ahorrando unas 2 horas de depuración manual por incidente.
            \end{highlights}
        \end{threecolentry}

    \medskip

    \section{Proyectos Destacados}

        \begin{threecolentry}{}{
            2026 - Presente
        }
            \textbf{MLB Lineup Optimizer — Sistema de Machine Learning End-to-End}\\
            \textit{Proyecto Personal}
            \begin{highlights}
                \item \textbf{Diseñé y desplegué un sistema de ML completo (ingesta $\rightarrow$ feature store $\rightarrow$ modelo $\rightarrow$ simulación $\rightarrow$ optimización $\rightarrow$ API)} sobre +2M de turnos al bate de 12 temporadas de MLB, con feature engineering anti-leakage y validación temporal estricta.
                \item \textbf{Entrené un modelo LightGBM multiclase calibrado (ECE de 0.002 en holdout)} cuyas probabilidades alimentan una simulación Monte Carlo (cadena de Markov base-outs, 100.000 juegos por partido) para estimar carreras esperadas y probabilidad de victoria.
                \item \textbf{Optimicé el orden de bateo con un algoritmo genético, logrando una mejora estadísticamente significativa de $\sim$+8.5 carreras/temporada} (IC90 vía bootstrap) frente al baseline, validada con un framework de backtesting propio.
                \item \textbf{Implementé prácticas MLOps de producción}: orquestación del pipeline con \textbf{Apache Airflow} sobre \textbf{Docker} (DAGs programados), MLflow tracking, gate de despliegue por control de drift, promoción champion/challenger y arquitectura de microservicios (FastAPI + dashboard React).
                \item \textbf{Repo: } \url{https://github.com/fernandogg2004/mlb-lineup-optimizer}
            \end{highlights}
        \end{threecolentry}

        \medskip

        \begin{threecolentry}{}{
            Ene 2026 - Jun 2026
        }
            \textbf{Modelos de Regresión No Lineales (Trabajo de Fin de Grado)}\\
            \textit{Universidad de Santiago de Compostela}
            \begin{highlights}
                \item Investigación y desarrollo teórico-práctico sobre el ajuste de funciones no lineales, metodologías de inferencia sobre parámetros y validación estricta de hipótesis estadísticas.
            \end{highlights}
        \end{threecolentry}

        \medskip

        \begin{threecolentry}{}{
        2025
        }
            \textbf{Modelización Estadística del Rendimiento Deportivo}\\
            \textit{Proyecto Académico}
            \begin{highlights}
                \item Desarrollo íntegro de un modelo estadístico para predecir y analizar el número de goles marcados por equipos de fútbol en base a variables de rendimiento en R.
                \item Mejora de más de un 14\% en $R^2$ respecto a baseline en un MSE de $\sim$0.5.
                \item \textbf{Repo: } \url{https://github.com/fernandogg2004/RendimientoDeportivo}
            \end{highlights}
        \end{threecolentry}

        \newpage

        \begin{threecolentry}{}{
        2024
        }
            \textbf{Modelado Analítico de Crecimiento Poblacional}\\
            \textit{Proyecto de Investigación Teórica}
            \begin{highlights}
                \item  Modelado matemático avanzado en sistema de EDOs y análisis de estabilidad de la evolución de una especie biológica ante un agente contaminante externo.
                \item \textbf{Repo: } \url{https://github.com/fernandogg2004/CrecimientoPoblacional}
            \end{highlights}
        \end{threecolentry}

    \medskip
   
    \section{Habilidades}

    \begin{threecolentry}{}{}
        \textbf{Competencias Técnicas}
        \begin{highlights}
            \item \textbf{Lenguajes de programación}: Python, R.
            \item \textbf{Librerías y Machine Learning:} Scikit-learn (sklearn), Pandas, NumPy, Polars, XGBoost, LightGBM.
            \item \textbf{Visualización:} Matplotlib, Seaborn, Plotly.
            \item \textbf{MLOps y APIs:} Docker, Apache Airflow (orquestación de pipelines), FastAPI, MLflow, simulación Monte Carlo y algoritmos genéticos, backtesting estadístico.
            \item \textbf{Software y Herramientas Analíticas:} Consumo de APIs REST con Requests / integración con JSON/Parquet, SQL (PostgreSQL, consultas analíticas con window functions), Jupyter + VS Code en entorno local o en la nube (Vertex AI Workbench), BigQuery, Git.
        \end{highlights}
    \end{threecolentry}

    \begin{threecolentry}{}{}
        \textbf{Idiomas y Estancia Internacional}
        \begin{highlights}
            \item \textbf{Inglés:} Competencia profesional completa (Nivel C1).
            \item \textbf{Español:} Nativo.
        \end{highlights}
    \end{threecolentry}



\end{document}